% ===========================================================================
\section{Discussion}
\label{sec:precise-alignment-discussion}
% ===========================================================================

This section interprets the quantitative results presented in \cref{sec:precise-alignment-results}, examining convergence behaviour, the effects of experimental variables (distance, angle), and the comparative strengths and limitations of the LiDAR and stereo modalities. A recurring theme is that neither modality dominates across all performance axes: LiDAR achieves higher convergence reliability through its multi-hypothesis perturbation strategy, whereas stereo produces lower inlier RMSE when it converges, reflecting its finer-grained depth resolution within a near-field window.

% ---------------------------------------------------------------------------
\subsection{Convergence Reliability}
\label{subsec:pa-discussion-convergence}
% ---------------------------------------------------------------------------

Convergence is the prerequisite for any downstream use of the alignment result. A pipeline that achieves low RMSE on converging trials but fails silently on a substantial fraction of conditions offers limited operational utility. In this regard, the LiDAR pipeline demonstrates a clear advantage: nine of ten conditions converged, with 311 of 315 hypotheses accepted when the non-converging D110A0 condition is excluded (98.7\% hypothesis-level acceptance). The stereo pipeline converged in seven of ten conditions (70\%), with the three failures occurring in distinct portions of the experimental design space.

The D110A0 condition represents a universal failure point: neither modality achieved registration at 110\,cm sensor-to-vine distance. The failure mechanisms, however, differ fundamentally between modalities. For LiDAR, the preprocessed source cloud contained 349\,663 points---approximately three times the 97\,000--133\,000 points observed in all other conditions. This anomalously large cloud indicates that the Patchwork++ ground segmentation algorithm~\cite{leePatchworkFastRobust2022} failed to remove ground-plane points at this distance, populating the source cloud with geometrically uninformative planar structure that obscures the vine trunk geometry and inflates point-to-plane residuals during GICP optimisation. For stereo, the failure is the converse: only 3\,484 raw points were generated (456 after processing), indicating that FoundationStereo~\cite{wenFoundationStereoZeroShot2025} cannot resolve sufficient geometric structure at 110\,cm given the stereo baseline and resolution of the ZED cameras. In both cases, the failure is attributable to the preprocessing stage rather than to limitations of the GICP algorithm itself---had the source clouds been correctly formed, registration may have succeeded.

The two stereo-only failures reveal different vulnerability modes. D40A$+$15 produced 87\,994 processed points---a density comparable to the successfully converging D40A$-$15 (87\,772 points) and D40A0 (83\,163 points)---yet achieved $F = 0.0$. Point density alone was not the limiting factor; rather, the combination of close range and positive approach angle appears to have placed the initial pose estimate in a region of the cost surface from which ICP diverged. Without any perturbation of the initial hypothesis, the stereo pipeline offers no recovery mechanism from such a local minimum. D80A$-$15 represents a different failure: only 23\,980 points survived preprocessing (compared to 42\,307 for D80A0 and 29\,981 for D80A$+$15), suggesting that the trunk segmentation stage was overly aggressive at this particular combination of distance and viewing angle, removing geometrically relevant structure and leaving insufficient correspondences for registration.

The architectural difference between the two pipelines---35 perturbation hypotheses for LiDAR versus a single hypothesis for stereo---accounts for a substantial portion of the convergence gap. The multi-hypothesis strategy functions as a form of random restart: even if the nominal initial pose lies near a local minimum or in a flat region of the cost surface, at least one of the 35 perturbed initialisations is likely to fall within the convergence basin of the global minimum. The D40A$+$15 case illustrates this directly---LiDAR achieved 35/35 accepted hypotheses under identical geometric conditions where stereo's single hypothesis diverged entirely. The multi-hypothesis approach thus provides robustness not through improved sensing, but through combinatorial coverage of the initialisation space.

The implication for operational deployment is that the stereo pipeline's lower convergence rate is partially addressable without changing the sensor modality. Introducing even a modest perturbation set (e.g., 5--7 hypotheses) to the stereo pipeline would provide robustness against the single-hypothesis vulnerability demonstrated by D40A$+$15, at the cost of a proportional increase in computation time. Alternatively, improving the accuracy of the initial pose estimate---through tighter coupling with the SLAM system or a coarse visual alignment step---would reduce the probability of the initial hypothesis falling outside the convergence basin. The convergence gap between modalities is thus partially architectural and partially physical, and the architectural component is readily addressable.

% ---------------------------------------------------------------------------
\subsection{Distance Effects on Registration Quality}
\label{subsec:pa-discussion-distance}
% ---------------------------------------------------------------------------

A counter-intuitive finding emerges from the distance analysis: the D80 group (80\,cm) achieves the lowest median inlier RMSE for both modalities---23.77\,mm for LiDAR and 12.45\,mm for stereo---outperforming both the closer D40 group (34.66\,mm LiDAR, 16.77\,mm stereo) and the D60 group (34.29\,mm LiDAR, 16.67\,mm stereo). An intuitive expectation is that closer sensing yields higher precision due to higher point density on the target structure; however, the data consistently indicate that moderate distance produces superior registration quality.

For the LiDAR modality, the geometric explanation centres on the spatial extent of structure captured within the sensor's field of view. At 40\,cm, the Livox Mid-360's field of view captures a relatively small, featureless cylindrical patch of the vine trunk. Smooth cylindrical geometry is well known to be problematic for ICP-class algorithms~\cite{rusinkiewiczEfficientVariantsICP2001}: point-to-plane correspondences along a cylinder are underconstrained in the tangential direction, leading to residual ambiguity that manifests as elevated RMSE. At 80\,cm, the wider field of view encompasses more of the vine's branching structure---lateral canes, cordons, and bifurcation points---providing geometrically distinctive features that constrain the registration more tightly. The IQR data support this interpretation: D80 conditions exhibit near-zero IQR in both fitness and RMSE (\cref{tab:pa-results-lidar-per-condition}), indicating that the solver converges to a well-defined minimum, whereas D40 conditions show slightly elevated IQR values consistent with residual geometric ambiguity.

The stereo modality exhibits the same distance-RMSE trend (D80 $<$ D60 $<$ D40), though the absolute RMSE values are approximately half those of LiDAR across all distance groups. The stereo pipeline applies a near-field depth clipping threshold ($z_\text{far} = 1.0$\,m), retaining only high-confidence depth estimates within the stereo baseline's optimal operating range. At D80, this clipping window captures the vine structure at a distance where the stereo disparity estimation is geometrically well-conditioned~\cite{wenFoundationStereoZeroShot2025}---the baseline-to-depth ratio produces sub-pixel disparity that translates to millimetre-level depth precision. At D40, the same clipping window may include points at the extreme edges of the depth range where stereo triangulation uncertainty increases, and the very close range may introduce depth-edge artefacts from the narrow stereo baseline. The net effect is a modest but consistent RMSE advantage for the D80 group.

The D110 condition represents a hard upper boundary for both sensing modalities, with qualitatively different failure mechanisms as discussed in \cref{subsec:pa-discussion-convergence}. This establishes an asymmetric operating envelope: registration degrades gracefully as distance decreases below 80\,cm (higher RMSE but maintained convergence), but fails abruptly as distance increases beyond approximately 100\,cm. The asymmetry has practical significance for path planning: approaching closer than the optimal distance is preferable to exceeding it, as over-approach produces suboptimal but usable alignments whereas under-approach produces no alignment at all.

The practical recommendation emerging from these results is an optimal operating range of 60--80\,cm for both modalities. Within this range, convergence is reliable ($\geq 97\%$ for LiDAR, 67--100\% for stereo) and RMSE is minimised. For context, prior work on LiDAR-based vine reconstruction by \citet{morenoOnGroundVineyardReconstruction2020} operated at comparable sensor-to-canopy distances (0.5--1.5\,m) in a mobile mapping configuration, though direct RMSE comparison is complicated by differences in registration targets (full canopy vs.\ isolated trunk). The sub-25\,mm LiDAR RMSE and sub-13\,mm stereo RMSE achieved at D80 place registration precision within the same order of magnitude as the target structures (5--15\,mm cane diameter), though they exceed the sub-centimetre goal stated in \cref{sec:research-objectives}.

% ---------------------------------------------------------------------------
\subsection{Approach Angle Effects}
\label{subsec:pa-discussion-angle}
% ---------------------------------------------------------------------------

The experimental design tested three approach angles ($-15^{\circ}$, $0^{\circ}$, $+15^{\circ}$) to assess whether oblique sensor-to-vine geometry degrades registration performance. For the LiDAR modality, the answer is negative within the tested range: both the A$-$15 and A$+$15 groups achieve 100\% hypothesis acceptance, with RMSE values (34.66\,mm and 34.28\,mm respectively) indistinguishable from the A0 group (34.22\,mm). The apparent reduction in A0 acceptance rate (72\%) is entirely attributable to the inclusion of D110A0, which failed for distance-related reasons unrelated to approach angle (\cref{tab:pa-results-lidar-by-angle}). The LiDAR pipeline is thus effectively invariant to $\pm 15^{\circ}$ approach angle variation, which is an important property for operational robustness given that a mobile platform navigating between vine rows cannot guarantee a perfectly perpendicular approach.

The stereo modality shows broadly similar convergence rates across angles (67\% for A$-$15 and A$+$15, 75\% for A0), but the individual failures reveal an asymmetry that merits interpretation. D40A$+$15 failed while D40A$-$15 succeeded under nominally symmetric angular conditions. This asymmetry suggests an interaction between approach angle and the specific vine geometry at close range: a $+15^{\circ}$ approach from 40\,cm may position the cameras such that the trunk is partially occluded by a lateral cane or the initial pose estimate is offset into a region where ICP diverges. The single-hypothesis vulnerability of the stereo pipeline amplifies such geometry-specific effects---with 35 hypotheses, LiDAR successfully registered D40A$+$15 on all trials, demonstrating that the underlying geometric content is adequate for registration and that stereo's failure is an initialisation problem rather than a sensing limitation.

The $\pm 15^{\circ}$ range tested here was chosen to represent a conservative estimate of approach angle variation for a mobile manipulator operating in vineyard rows. In practice, uneven or rocky terrain may cause the robot chassis to deviate from a perpendicular orientation relative to the vine, and imprecise stopping (slightly early or late along the row) introduces an angular offset between the sensor boresight and the vine centre. The $\pm 15^{\circ}$ envelope is intended to encompass these worst-case geometric deviations. Within this envelope, the results confirm that neither pipeline requires angular pre-alignment or view-planning prior to executing precise alignment. However, larger angular deviations ($\pm 30^{\circ}$ or beyond) remain untested and may introduce challenges---particularly for stereo, where the overlap between the sensed point cloud and the reference model decreases as the viewing direction diverges from the reference model's construction viewpoints. Characterising performance at larger angles would require additional experimental conditions beyond the scope of this evaluation.

% ---------------------------------------------------------------------------
\subsection{Perturbation Strategy Effectiveness}
\label{subsec:pa-discussion-perturbation}
% ---------------------------------------------------------------------------

The multi-hypothesis perturbation strategy generates 35 initial pose hypotheses per trial, distributed across five groups (Orig, A--E) that sample translational and rotational offsets around the SLAM-provided initial pose. The per-group acceptance rates are uniform: 88\% (Group~C) to 90\% (Orig, A, B), with median fitness and RMSE values statistically indistinguishable across groups (\cref{tab:pa-results-groups}). No perturbation group is systematically more or less likely to converge, indicating that the perturbation magnitudes are appropriately sized relative to the convergence basin width---large enough to escape local minima near the nominal pose, but not so large as to place hypotheses entirely outside the basin of attraction.

The within-trial convergence analysis (\cref{tab:pa-results-lidar-sigma}) reveals that the perturbation strategy provides \emph{redundancy} rather than \emph{diversity}. For the majority of conditions, accepted hypotheses converge to effectively the same final pose: D80A0 achieves near-zero standard deviation across all transformation components ($\sigma < 0.1$\,mm translation, $< 0.01^{\circ}$ rotation), and D40A0, D60A$-$15, and D60A0 show similarly tight clustering. This behaviour indicates that the cost surface in the vicinity of the correct alignment contains a single dominant basin, and the perturbation strategy's value lies not in discovering alternative solutions but in ensuring that at least one hypothesis starts within that basin. The redundancy provides a probabilistic guarantee: if any single hypothesis has probability $p$ of converging, 35 independent hypotheses achieve convergence probability $1 - (1-p)^{35}$, which exceeds 99\% for any per-hypothesis convergence probability above 13\%.

The exceptions to single-basin convergence are informative. D80A$-$15 and D80A$+$15 exhibit elevated roll standard deviation ($\sigma_\phi > 6^{\circ}$), and D40A$-$15 and D40A$+$15 show moderate roll variability ($\sigma_\phi \approx 4^{\circ}$). These conditions share a common geometric characteristic: the vine trunk presents an approximately cylindrical cross-section when viewed from an oblique angle. Cylindrical geometry is rotationally symmetric about its axis, creating a degenerate direction in the ICP cost surface along which the solver cannot distinguish between roll orientations. Some perturbation hypotheses settle at different roll values within this degenerate manifold, producing the observed multi-modal roll distribution. This roll ambiguity does not affect the translational components ($\sigma_{t_x}$, $\sigma_{t_y}$, $\sigma_{t_z}$ remain low for these conditions), confirming that the positional alignment is well-determined even when the rotational alignment around the trunk axis is not.

For operational deployment, the redundancy interpretation suggests that the number of hypotheses could be substantially reduced without sacrificing reliability. If 35 hypotheses converge to the same solution in 8 of 9 conditions, then 7--10 hypotheses would provide equivalent coverage with a threefold to fivefold reduction in computation time. The primary exception---D80A0, where only 32 of 35 hypotheses were accepted---still shows a 91\% per-hypothesis acceptance rate, for which 10 hypotheses would yield $>$99.99\% trial-level convergence probability. Such a reduction would bring the LiDAR pipeline's timing from $\sim$46\,s down to $\sim$12--15\,s, approaching the stereo pipeline's latency and substantially improving real-time feasibility without requiring algorithmic parallelisation.

% ---------------------------------------------------------------------------
\subsection{Cross-Modal Trade-offs}
\label{subsec:pa-discussion-cross-modal}
% ---------------------------------------------------------------------------

The matched-condition comparison (\cref{tab:pa-results-cross-modal}) summarises the central finding of this evaluation: the LiDAR and stereo modalities occupy complementary positions in the reliability--precision--speed trade-off space. Neither modality dominates the other across all axes of performance, and the choice between them (or the strategy for combining them) depends on the operational priority of the downstream application.

On the seven conditions where both pipelines converged, stereo achieves a mean inlier RMSE of 15.51\,mm compared to 31.40\,mm for LiDAR---a factor-of-two precision advantage. This difference is not primarily attributable to the ICP algorithm (both pipelines use the same GICP solver~\cite{segalGeneralizedICP2009} with identical convergence parameters) but to the quality of the input point clouds. The stereo pipeline's near-field depth clipping ($z_\text{far} = 1.0$\,m) retains only points within the optimal operating range of the stereo triangulation geometry, producing a compact, high-precision cloud concentrated on the vine structure. The LiDAR pipeline, by contrast, operates on an accumulated cloud that includes points from a wider angular range and greater depth variation, incorporating returns with higher geometric uncertainty from oblique incidence angles and surface discontinuities. Importantly, RMSE as reported here measures the \emph{internal consistency} of inlier correspondences---the mean distance between matched source-target point pairs after alignment---rather than absolute positional accuracy relative to a ground-truth pose. A low RMSE indicates tight geometric agreement between the registered clouds, which is a necessary condition for pose accuracy but not a sufficient one.

The reliability advantage of LiDAR is equally clear: 90\% condition-level convergence versus 70\% for stereo. LiDAR converged in two conditions where stereo failed (D40A$+$15 and D80A$-$15), while no condition that failed for LiDAR succeeded for stereo. This asymmetry is partially explained by the multi-hypothesis architecture discussed in \cref{subsec:pa-discussion-perturbation}, and partially by the sensing physics. LiDAR is an active sensor that generates its own illumination at 905\,nm, making it insensitive to ambient lighting conditions and surface texture. Stereo, as a passive modality, depends on photometric contrast for disparity estimation; in a laboratory environment this is not limiting, but in field conditions where lighting varies and surfaces may be wet or uniformly coloured, the reliability gap could widen further. The result that LiDAR succeeds everywhere stereo succeeds (but not vice versa) suggests that stereo failures are initialisation or preprocessing problems rather than fundamental geometric insufficiency.

The timing differential is substantial. The stereo pipeline completes in $13.69 \pm 2.82$\,s with no preceding accumulation interval, whereas the LiDAR pipeline requires $46.0 \pm 3.8$\,s of processing plus a 25\,s scan accumulation phase (though accumulation can overlap with prior-vine pruning). In a pruning operation where the manipulator requires 30--60\,s to plan and execute a cutting trajectory, the stereo pipeline's latency can be entirely absorbed within the approach and planning phases, whereas the LiDAR pipeline's current sequential implementation would place alignment on the critical path. The approximately threefold speed advantage of stereo is primarily structural: stereo executes a single ICP run versus LiDAR's 35 sequential hypotheses, and processes fewer points (68\,000--92\,000 vs.\ 97\,000--133\,000).

The complementary profiles suggest a natural deployment architecture: a cascaded strategy in which stereo alignment is attempted first, and LiDAR is invoked as a fallback only when the stereo fitness score indicates non-convergence. Under such an architecture, the majority of alignment events would complete in $\sim$14\,s with stereo precision ($\sim$15\,mm RMSE), while the subset of conditions that challenge stereo (estimated at 30\% based on this evaluation) would be rescued by LiDAR at the cost of additional latency and somewhat reduced precision. The cascaded approach would achieve higher aggregate convergence than either modality alone---potentially exceeding the 90\% and 70\% rates observed for LiDAR and stereo individually within the 40--80\,cm operating range---while preserving stereo's speed advantage for the common case. An alternative fusion strategy, such as confidence-weighted averaging of the two alignment results when both converge, is also conceivable but would require validation against external ground truth to confirm that combining two internally consistent but potentially differently biased estimates improves absolute accuracy.

The cross-modal comparison also validates the two-stage localisation architecture proposed in this thesis. The ICP corrections reported in \cref{tab:pa-results-lidar-icp-correction,tab:pa-results-stereo-transforms} range from 46.6\,mm to 165.3\,mm in translational displacement and from $0.13^{\circ}$ to $11.5^{\circ}$ in rotational correction. These magnitudes are consistent with the SLAM system's (GLIM-GNSS) absolute trajectory error of 71--174\,mm reported in \cref{chap:SLAM}, confirming that the coarse localisation delivered by SLAM places the robot within ICP's convergence basin. The perturbation strategy samples $\pm 50$\,mm translation and $\pm 10^{\circ}$ rotation around the SLAM pose, and achieves 98.7\% hypothesis-level convergence---evidence that the SLAM accuracy is sufficient for the precise-alignment stage to function reliably. The two-stage architecture thus combines capabilities that neither stage provides independently: SLAM provides global consistency and centimetre-level initial pose, while ICP refines this to millimetre-level alignment against the reference model.

% ---------------------------------------------------------------------------
\subsection{Pipeline Timing and Real-Time Feasibility}
\label{subsec:pa-discussion-timing}
% ---------------------------------------------------------------------------

The concept of ``real-time'' in the context of robotic vine pruning is defined not by absolute latency thresholds but by the pruning cycle's critical path. For the purposes of this analysis, a representative per-vine cycle is assumed to consist of: approach and stop (5--10\,s), precise alignment, motion planning and trajectory generation (5--15\,s), and manipulator execution (30--60\,s). Alignment is on the critical path only if it cannot overlap with other phases and its latency exceeds the time budget allocated between the robot's arrival at the vine and the initiation of manipulator motion. Under this framing, an alignment latency of 10--15\,s is effectively real-time (absorbed within approach and planning), whereas latency exceeding 30\,s extends the per-vine cycle and reduces throughput.

The stereo pipeline, with a mean total latency of $13.69 \pm 2.82$\,s, already satisfies this real-time criterion. Its timing decomposition reveals that the dominant components are ICP alignment ($7.16 \pm 2.37$\,s) and dual-camera combination ($4.41 \pm 0.58$\,s), with GPU-accelerated depth estimation via FoundationStereo contributing only 1.2\,s. The combination stage---which merges point clouds from two stereo cameras into a unified source cloud---is a software implementation detail rather than a fundamental algorithmic limitation; optimisation through more efficient data structures or pre-allocation could reduce this component. The ICP latency is determined by cloud size and convergence speed, both of which are already favourable for the stereo pipeline's relatively compact clouds (68\,000--92\,000 processed points). No architectural changes are required for the stereo pipeline to operate within the pruning cycle's time budget.

The LiDAR pipeline's current latency of $46.0 \pm 3.8$\,s (plus 25\,s scan accumulation) substantially exceeds the real-time threshold. However, the timing structure reveals that this latency is an artefact of sequential execution rather than a fundamental computational limit. The 35 perturbation hypotheses are evaluated independently and sequentially, each requiring approximately 0.96\,s of ICP computation. These hypotheses are \emph{embarrassingly parallel}: they share the same source and target clouds, differ only in their initial transformation, and produce independent results. On modern multi-core hardware or with GPU-accelerated ICP implementations such as those provided by GLIM~\cite{koideGLIM3DRangeInertial2024}, all 35 hypotheses could execute concurrently, reducing the registration phase from $\sim$34\,s to a single ICP iteration of $\sim$1\,s. Combined with the 0.71\,s preprocessing overhead, the parallelised LiDAR pipeline would achieve $\sim$2\,s active computation time---faster than the stereo pipeline. The 25\,s accumulation interval can be overlapped with the preceding vine's pruning execution, effectively removing it from the critical path. With parallelisation and accumulation overlap, the LiDAR pipeline's effective per-vine latency would approach 2--3\,s, well within real-time constraints. Furthermore, the hypothesis reduction discussed in \cref{subsec:pa-discussion-perturbation}---from 35 to 7--10 hypotheses---would achieve a similar latency improvement ($\sim$10--12\,s) without requiring parallel hardware, representing a path to real-time operation that requires minimal engineering effort.

The D110A0 timing anomaly (63.5\,s for LiDAR, 9.3\,s for stereo) illustrates how preprocessing failures cascade into timing. The anomalously large LiDAR source cloud (349\,663 points) inflates both preprocessing time (2.29\,s vs.\ $\sim$0.71\,s typical) and per-hypothesis ICP time (1.29\,s vs.\ $\sim$0.96\,s typical). For stereo, the very small cloud (456 points) paradoxically reduces ICP time because the solver terminates quickly on a degenerate input. In both cases, the timing anomaly serves as a secondary diagnostic of preprocessing failure---an unexpectedly fast or slow pipeline execution could trigger fallback behaviour in a production system, even before examining the fitness score.

% ---------------------------------------------------------------------------
\subsection{Limitations}
\label{subsec:pa-discussion-limitations}
% ---------------------------------------------------------------------------

The most significant limitation of this evaluation is the absence of an external six-degree-of-freedom ground-truth measurement. All reported metrics---fitness score and inlier RMSE---are \emph{internal consistency} measures: they quantify how well the source and target point clouds agree geometrically after registration, but cannot determine whether the recovered pose is correct in an absolute sense. A registration that achieves high fitness and low RMSE could, in principle, converge to a self-consistent but incorrect local minimum, particularly if both clouds share systematic biases (e.g., both derived from the same sensor platform with an uncalibrated extrinsic offset). The visual alignment overlays (\cref{fig:pa-lidar-alignment-examples,fig:pa-stereo-converged-overlay}) provide qualitative evidence that the registrations are geometrically plausible, but qualitative assessment cannot substitute for quantitative validation against an independent pose reference such as a motion capture system or precision-surveyed fiducial markers. Future work should incorporate such ground truth to establish absolute pose accuracy bounds and determine whether the RMSE values reported here correspond to comparable levels of absolute error.

The evaluation was conducted on a single dormant vine specimen in a controlled laboratory environment. While this approach eliminates confounding factors (wind, variable lighting, inter-vine occlusion) and permits precise control of experimental variables, it leaves open the question of generalisation. Field conditions introduce multiple sources of variability absent from the laboratory: wind-induced motion of canes and foliage (particularly relevant for non-dormant canopy), varying vine structure across cultivars, training systems, and ages, canopy interference from adjacent rows, variable ambient lighting affecting stereo depth estimation, and surface contamination (mud, water droplets) affecting LiDAR reflectance. Each of these factors could independently degrade registration performance, and their combined effect in field conditions is unknown. The single-vine design also means that the results reflect the specific geometry of the test specimen---a spur-pruned Sauvignon Blanc vine on a vertical shoot positioning trellis---and may not transfer directly to other pruning architectures or cultivars with substantially different growth habits.

The experimental design employs a $3 \times 3 + 1$ factorial with a single observation per condition ($n = 1$). This design is sufficient to characterise broad performance trends (distance effects, modality differences) but insufficient for rigorous statistical inference. Without replicates, it is impossible to distinguish genuine condition effects from trial-specific noise---for example, the D40A$+$15 stereo failure could reflect a systematic vulnerability at that condition or a stochastic event that would not recur on a second trial. The absence of repeated measurements also precludes assessment of measurement repeatability: the extent to which the same condition produces the same alignment result across multiple independent trials. For a pruning system that must reliably position a cutting tool within millimetres of a target, repeatability is as important as accuracy for operational deployment, and this evaluation provides no information on that axis.

The condition-specific failures (D110A0 for both modalities, D40A$+$15 and D80A$-$15 for stereo) may be partially specific to the test vine's geometry and the chosen preprocessing parameters. The D110A0 failure, while clearly attributable to preprocessing limitations (ground segmentation failure for LiDAR, insufficient depth resolution for stereo), occurred at a distance threshold that may shift for different vine architectures or sensor mounting heights. Similarly, the D40A$+$15 stereo failure may reflect a geometric interaction between the particular arrangement of canes on this vine and the camera viewpoint at that condition, rather than a general vulnerability of the stereo pipeline at 40\,cm with $+15^{\circ}$ approach angle. Characterising the robustness of these failure boundaries would require testing across multiple vine specimens, which is left as future work.

Finally, the preprocessing pipeline represents a confounding factor in attributing performance characteristics to the sensing modality versus the algorithmic choices. Failures attributed to ``the LiDAR pipeline'' conflate the Livox Mid-360's sensing capabilities with the specific Patchwork++ parameterisation, the statistical outlier rejection thresholds, and the voxel downsampling resolution; similarly, ``stereo pipeline'' failures conflate FoundationStereo's depth estimation with the DBSCAN clustering parameters, morphological closing kernel size, and the $z_\text{far}$ clipping threshold. A different parameterisation of these preprocessing stages might shift the convergence boundaries or alter the RMSE characteristics. Disentangling the contributions of sensing physics from algorithmic preprocessing choices would require ablation studies that systematically vary preprocessing parameters while holding the sensing modality constant---an analysis beyond the scope of this evaluation but necessary for informed pipeline optimisation in future deployments.
